\documentclass[11pt]{article}

\usepackage[margin=1in]{geometry}
\usepackage{amsmath,amssymb,amsfonts,amsthm}
\usepackage{booktabs}
\usepackage{array}
\usepackage{xcolor}
\usepackage{hyperref}
\usepackage{enumitem}
\usepackage{caption}
\usepackage{graphicx}

\hypersetup{
  colorlinks=true,
  linkcolor=blue!50!black,
  citecolor=blue!50!black,
  urlcolor=blue!50!black
}

\newtheorem{theorem}{Theorem}
\newtheorem{lemma}{Lemma}
\newtheorem{definition}{Definition}
\newtheorem{corollary}{Corollary}

\newcommand{\Fp}{\mathbb{F}_p}
\newcommand{\MI}{M_{I}}
\newcommand{\Mfull}{M_{\mathrm{full}}}

\title{\textbf{Beyond MDS: A Unified Generation-Order and Structural\\
Criterion for Poseidon-Family Linear Layers}}
\author{}
\date{November 2026}

\begin{document}
\maketitle

\begin{abstract}
The partial rounds of the Poseidon hash family apply a nonlinear S-box to a
single state coordinate and rely on a linear layer $M$ to diffuse that
nonlinearity across the remaining coordinates. The folklore security
argument is that $M$ must be MDS: a non-MDS $M$ admits an explicit invariant
subspace trail that collapses an arbitrary number of partial rounds to a
linear map, which we reproduce as a baseline result. We then show that MDS
is \emph{necessary but not sufficient}: three independent 2026 papers
(Vitto; Merz and Rodr\'iguez Garc\'ia; Jo) each break properly-MDS,
subspace-trail-clean Poseidon instances through three distinct mechanisms
--- finite-trail full-round absorption, tensor-structured round skipping,
and a full-round compression collision from an adaptively chosen matrix,
respectively. We give a single generation-order countermeasure
(cryptographically binding the round constants to a canonical encoding of
the linear layer, so no matrix can be chosen after the constants it will
interact with are already fixed) and validate it against all three attacks,
including a from-scratch replication of Jo's attack and an empirical
immunity test of our own binding against it --- to our knowledge the first
treatment of that attack beyond a one-sentence mention in the literature.
We combine this with a structural criterion (generalizing both the
GRS21 invariant-subspace check and the tensor-eigenvector check) and apply
both to a lightweight KoalaBear instantiation, cross-validated against a
real \texttt{circom}/Groth16 compile, the production \texttt{circomlib}
reference implementation, and an exhaustive eight-point adversarial
falsification battery. We report real R1CS constraint counts and position
the result against Poseidon, Poseidon2, Anemoi, and Polocolo.
\end{abstract}

\tableofcontents

% =========================================================================
\section{Introduction}
% =========================================================================

Poseidon~\cite{gkr21} is an arithmetization-oriented permutation built on
the HADES strategy~\cite{hades20}: \emph{full rounds} apply a power S-box
$S(x) = x^{\alpha}$ to every coordinate of the state, while \emph{partial
rounds} apply it to a single coordinate only, leaving the rest linear.
Every round (full or partial) is followed by the addition of round
constants and multiplication by a fixed linear layer $M$. The efficiency
of Poseidon on arithmetic circuits rests entirely on partial rounds being
cheap (one nonlinear constraint per round instead of $t$), so the
diffusion properties of $M$ carry the full burden of mixing the single
active coordinate back into the rest of the state.

\paragraph{The folklore criterion.} The standard requirement on $M$ is
that it be \emph{maximum distance separable} (MDS). We reproduce, as a
self-contained baseline (\S\ref{sec:invariant-subspace}), why this is
\emph{necessary}: a non-MDS matrix admits a nonzero difference that never
reaches the active coordinate, so a chosen-difference adversary can
predict the permutation's output without ever engaging the S-box.

\paragraph{MDS is not sufficient.} The contribution of this paper begins
from a 2026 observation that the folklore criterion is incomplete. Three
independent papers --- developed in parallel, none derivative of the
others~\cite{vitto2026,merzgarcia2026,jo2026} --- each construct a
properly MDS linear layer, passing every published subspace-trail check,
that nonetheless breaks the differential or collision security of a
standard Poseidon instance:

\begin{itemize}[leftmargin=1.4em]
  \item \textbf{Vitto (Slipway, 2026/1579)}~\cite{vitto2026}: every linear
  layer unavoidably admits \emph{finite} subspace trails. If the matrix is
  chosen \emph{after} the round constants are fixed, an adversary can
  steer a constrained input family through the initial full rounds so
  that it lands inside that trail, collapsing the effective algebraic
  degree far below the naive bound.
  \item \textbf{Merz and Rodr\'iguez Garc\'ia (Skipping Class,
  2026/306)}~\cite{merzgarcia2026}: Poseidon2(b)'s~\cite{gks23} efficient
  Kronecker-structured external matrix $M_\varepsilon = P_{t/4}\otimes M_4$
  has an unconstrained $(q,-q,0,\dots,0)$ eigenvector subspace that
  survives full rounds at zero algebraic-degree cost, independently of the
  (already-vetted) internal matrix.
  \item \textbf{Jo (Midpoint Reset, 2026/1760)}~\cite{jo2026}: if the
  round constants are fixed before the matrix is chosen, an adversary can
  program an \emph{exact} nonlinear two-execution trajectory through every
  round via a linear interpolation $M = VU^{-1}$, yielding a genuine
  full-round \emph{compression collision} --- not merely a distinguisher
  --- in a matrix that is MDS and passes every subspace-trail check in the
  literature.
\end{itemize}

Each of the three papers names its own ad hoc countermeasure for the
mechanism it exhibits. None states or validates a criterion that covers
more than one mechanism, and (as we confirm in
\S\ref{sec:sok} against a full read of the most comprehensive available
survey~\cite{sok2026}) no existing source analyzes Jo's attack beyond a
one-sentence mention, let alone proposes and tests a defense against it.

\paragraph{Contributions.}
\begin{enumerate}[leftmargin=1.4em]
  \item A \textbf{unified structural criterion} (\S\ref{sec:criterion})
  generalizing both the GRS21 invariant-subspace
  check~\cite{grs21} and the tensor-eigenvector check implicit in
  Merz and Rodr\'iguez Garc\'ia's attack, implemented as two
  field-generic algorithms that together flag every unconstrained
  invariant subspace surviving into or through the nonlinear rounds.
  \item A \textbf{single generation-order countermeasure} ---
  cryptographically binding round constants to a canonical encoding of
  \emph{both} the partial- and full-round matrices before either is
  used --- and its \textbf{empirical validation against all three attacks},
  including a from-scratch replication of Jo's construction and a direct
  test of whether our binding defeats it (\S\ref{sec:midpoint-reset}),
  which we believe is the first such validation reported anywhere.
  \item A \textbf{lightweight instantiation} at KoalaBear
  ($t=16,\alpha=3,R_F=8,R_P=20$), combining a NUMS partial-round matrix
  with a tensor-structured, safe-order full-round matrix, cross-validated
  against a real \texttt{circom}/Groth16 compile and the production
  \texttt{circomlib} reference, and scored against an
  \textbf{eight-point exhaustive falsification battery}
  (\S\ref{sec:battery}).
  \item An honest account of the design levers that \emph{did not} help
  (S-box exponent reduction, S-box-lane rotation, two tricks imported from
  GMiMCHash2, and a naive lookup-hybrid model), reported as negative
  results rather than omitted (\S\ref{sec:levers}).
\end{enumerate}

% =========================================================================
\section{Preliminaries}
% =========================================================================

\subsection{Poseidon1 round structure}

Let $\Fp$ be a prime field, $t$ the state width, $\alpha$ the S-box
exponent with $\gcd(\alpha, p-1) = 1$. A state is a vector
$x = (x_0,\dots,x_{t-1}) \in \Fp^t$. A \emph{full round} applies

\[
  S_F(z) = z^{\circ \alpha} \quad\text{(componentwise)},
\]

while a \emph{partial round} applies

\[
  S_P(z) = (z_0^{\alpha}, z_1, \dots, z_{t-1})^{T}.
\]

In this paper's own convention (matching the repository's existing code
and the original CLAUDE.md problem statement), a round is
\[
  R_r(x) = M \cdot \big(S_r(x) + c^{(r)}\big),
\]
i.e. S-box, then add round constants, then apply the linear layer $M$.
(Jo's paper~\cite{jo2026} instead uses $R_r(x) = M\, S_r\big(x + c^{(r)}\big)$,
constants added \emph{before} the S-box; \S\ref{sec:midpoint-reset}'s
replication module implements that convention directly rather than
reconciling the two, since Jo's Lemma is proved for that exact ordering.)
The full permutation applies $R_F/2$ full rounds, then $R_P$ partial
rounds, then $R_F/2$ more full rounds.

\subsection{The MDS / branch-number criterion}

\begin{definition}[Branch number]
For $M \in \Fp^{t\times t}$,
\[
  B(M) \;=\; \min_{\Delta \neq 0} \big(\mathrm{wt}(\Delta) + \mathrm{wt}(M\Delta)\big).
\]
$M$ is \emph{MDS} iff $B(M) \geq t+1$ (equivalently, every square
submatrix of $M$ is nonsingular).
\end{definition}

\subsection{Why non-MDS is unsafe: an invariant subspace trail}
\label{sec:invariant-subspace}

We restate the project's original baseline result as a self-contained
lemma, since the rest of the paper depends on the reader accepting that
MDS is at least \emph{necessary}.

\begin{lemma}[Invariant subspace trail under a non-MDS matrix]
\label{lem:invariant-subspace}
Let $t=3$ and let $M$ be non-MDS, so there exists a nonzero
$\Delta = (0,0,\delta)$ with $M\Delta \in V := \{(0,0,x) : x \in \Fp\}$.
Then $V$ is invariant under the full sequence of $R_P$ partial rounds:
for any two states $A,B$ with $A-B \in V$, the S-box never activates on
their difference, and the entire partial-round sequence collapses to the
linear map $\Delta \mapsto M_{22}^{R_P}\,\Delta$ restricted to $V$, where
$M_{22}$ is $M$'s bottom-right entry.
\end{lemma}

\begin{proof}
Write $A_0 = B_0$ (since $\Delta_0 = 0$), hence $A_0^{\alpha} = B_0^{\alpha}$
and the post-S-box difference is still $(0,0,\delta)$. Round-constant
addition is affine and preserves the difference. By hypothesis
$M\cdot(0,0,\delta) \in V$, so the difference re-enters $V$ at the start
of the next round. Induction over $R_P$ rounds gives the claim.
\end{proof}

\paragraph{Concrete confirmation.} For $M = I$ (the identity, a non-MDS
matrix with $B(I)=2 < t+1=4$), $M\cdot(0,0,\delta) = (0,0,\delta)$
trivially, and an adversary with any difference confined to $V$ can
predict the permutation's output difference with zero oracle queries: a
simulated diffusion trace over $6$ partial rounds at $\alpha=5$ shows the
diffusion coefficient $D(R)$ frozen at $1/3$ indefinitely under $M=I$,
versus reaching full diffusion ($D(R)=1$) after a single round under an
MDS benchmark matrix (Table~\ref{tab:diffusion}).

\begin{table}[h!]
\centering
\caption{Diffusion coefficient $D(R)$ over $6$ partial rounds, BN254,
$\alpha=5$, input difference $\Delta=(0,0,4)$.}
\label{tab:diffusion}
\begin{tabular}{@{}lcc@{}}
\toprule
Round & Non-MDS ($M=I$) & MDS benchmark \\
\midrule
0 & $1/3$ & $1/3$ \\
1 & $1/3$ & $3/3$ \\
2--6 & $1/3$ & $3/3$ \\
\bottomrule
\end{tabular}
\end{table}

This establishes necessity. The remainder of the paper is about why it is
not sufficient.

% =========================================================================
\section{Three Independent Attacks on Properly-MDS Instances}
\label{sec:threats}
% =========================================================================

Table~\ref{tab:threats} summarizes the three attacks this paper addresses,
alongside two attack classes that, as we discuss in
\S\ref{sec:sok}, remain genuinely open.

\begin{table}[h!]
\centering
\caption{Threat model. Items 1--3 are addressed by this paper
(\S\ref{sec:criterion}--\S\ref{sec:midpoint-reset}); items 4--5 remain
future work (\S\ref{sec:discussion}).}
\label{tab:threats}
\small
\begin{tabular}{@{}p{0.12\linewidth}p{0.4\linewidth}p{0.4\linewidth}@{}}
\toprule
\textbf{Attack} & \textbf{Mechanism} & \textbf{Our status} \\
\midrule
Infinite invariant subspace trail &
$B(M) < t+1$ lets a coordinate subspace stay invariant forever &
Closed: MDS + GRS21 irreducibility (Algorithm~1, \S\ref{sec:criterion}) \\
\addlinespace
Finite trail + full-round absorption (Vitto~\cite{vitto2026}) &
Finite trails are unavoidable; adaptive matrix choice steers a family into one via the full-round prefix &
Closed: binding (\S\ref{sec:binding}) + Algorithm~2's finite-trail profile \\
\addlinespace
Tensor round-skipping (Merz \& Garc\'ia~\cite{merzgarcia2026}) &
$M_\varepsilon = P_{t/4}\otimes M_4$ has an unconstrained $(q,-q,0,\dots)$ eigenspace surviving full rounds &
Closed: swapped factor order $M_4\otimes P_{t/4}$ (\S\ref{sec:tensor}) \\
\addlinespace
Adaptive-matrix collision (Jo~\cite{jo2026}) &
$M=VU^{-1}$ interpolation against fixed constants programs an exact collision trajectory; $M$ is MDS and passes every trail check &
Closed: binding, replicated and empirically tested (\S\ref{sec:midpoint-reset}) \\
\addlinespace
S-box skipping~\cite{sboxskip2026} (Bhati, Tariq, Ashur, 2026/1692) &
Generalized subspace-restriction gadget, claimed independent of $M$, constants, $\alpha$, field &
\textbf{Open} --- not analyzed in depth anywhere, including the SoK (\S\ref{sec:sok}) \\
\addlinespace
Bivariate resultants against Poseidon~\cite{resultant2026} (2026/1905) &
Resultant-based CICO solving, published against Type-II primitives, not Poseidon's own linear layer &
\textbf{Open} for Poseidon specifically \\
\bottomrule
\end{tabular}
\end{table}

% =========================================================================
\section{A Unified Structural Criterion}
\label{sec:criterion}
% =========================================================================

Vitto's and Merz and Rodr\'iguez Garc\'ia's mechanisms differ in which
matrix ($\MI$ vs.\ $\Mfull$) and which round segment (partial vs.\ full) an
unconstrained subspace survives through, but both instantiate the same
underlying question:

\begin{quote}
\emph{For a candidate $(\MI, \Mfull)$, does any proper subspace of either
matrix's eigenstructure survive into the nonlinear rounds without being
constrained by the generation procedure?}
\end{quote}

We implement this as two complementary, field-generic algorithms.

\paragraph{Algorithm 1 (infinite invariant subspace in $\MI$).}
Generalizes GRS21~\cite{grs21}: factor the characteristic polynomial
$\chi(\MI)$ over $\Fp$, and for each irreducible factor $f_i$, compute
$\ker f_i(\MI)$ via the minimal/characteristic-polynomial decomposition.
A fast path checks irreducibility of $\chi(\MI^{\,i})$ for
$i=1,\dots,4(t-1)$ (GRS21's sufficient-but-not-necessary shortcut): when it
holds, no nontrivial invariant subspace exists at any dimension, without
running the general decomposition.

\paragraph{Algorithm 2 (finite-trail profile and full-round scan).}
Two complementary pieces:
\begin{itemize}[leftmargin=1.4em]
  \item \emph{Finite-trail dimension profile.} $K_r(M) =
  \bigcap_{j=0}^{r-1}\ker\!\big(e_0^{T}M^{j}\big)$, the subspace of
  directions whose active-coordinate contribution stays $X$-independent
  for $r$ rounds (Vitto's Eq.~2). We verify $\dim K_r(\MI) = t - r$
  (``maximal partial-round rank growth'') as a baseline that confirms a
  candidate's finite trail is no longer than the unavoidable minimum.
  \item \emph{Full-round invariant-subspace scan.} Generalizes the
  tensor-eigenvector finding of~\cite{merzgarcia2026} to \emph{any} full-round
  matrix, with no Kronecker structure assumed: for every irreducible
  factor of $\chi(\Mfull)$, compute the corresponding invariant subspace
  and flag it when its basis vectors exhibit a contiguous zero-run of
  length $\geq t/4$ (the threshold Merz and Rodr\'iguez Garc\'ia's own
  CICO-$k$ analysis uses). This is a \emph{static, necessary-looking but
  not sufficient} condition: \S\ref{sec:battery} documents a confirmed
  false positive at KoalaBear scale, resolved only by pairing the static
  flag with a dynamic one-round nonlinear survival check.
\end{itemize}

\paragraph{What this criterion cannot see.} Both algorithms reason about
static properties of $M$ in isolation. Jo's attack constructs an $M$ with
\emph{no} exploitable structural property --- it is MDS, and both
algorithms above report it secure --- because the vulnerability is a
\emph{relation between $M$ and the already-fixed constants}, not a
property of $M$ alone. No static scan, however sophisticated, can close
this gap; it requires controlling the \emph{order} in which the matrix and
the constants are generated.

% =========================================================================
\section{The Generation-Order Countermeasure}
\label{sec:binding}
% =========================================================================

\subsection{Binding constants to the matrix}

Both Vitto's and Jo's attacks (and the generic kleptographic-backdoor
class~\cite{backdoor2026}) require the adversary to choose, or influence
the choice of, the linear layer \emph{after} the round constants it will
interact with are fixed. We remove this degree of freedom by construction:

\[
  T = \texttt{seed} \,\|\, \mathrm{enc}(p,t,\alpha,R_F,R_P) \,\|\, \mathrm{enc}(\MI) \,\|\, \mathrm{enc}(\Mfull),
  \qquad
  c^{(r)}_i = \mathrm{HashToField}\big(T \,\|\, \mathrm{enc}(r,i)\big).
\]

Both matrices are generated first, non-adaptively, via a
nothing-up-my-sleeve (NUMS) procedure; the round constants are then
derived deterministically from a transcript that already contains both
matrices. No constant can be fixed independently of either matrix, and no
matrix can be chosen as a function of constants that do not yet exist.
An earlier version of this construction bound only $\MI$; we found and
fixed this gap in a Phase-3 adversarial review, since Vitto's own attack
targets $\Mfull$ (the matrix governing the full-round absorption), not
$\MI$.

\subsection{Tensor-structured $\Mfull$: the safe factor order}
\label{sec:tensor}

For the Poseidon2(b)-style efficiency win of a Kronecker-structured
full-round matrix, $M_\varepsilon = P_{t/4}\otimes M_4$ admits the
$(q,-q,0,\dots,0)$ eigenvector subspace exploited
in~\cite{merzgarcia2026}; swapping the factor order to
$\bar M_\varepsilon = M_4 \otimes P_{t/4}$ is identical in multiplicative
cost and destroys the exploitable shape --- confirmed both by a static
scan (Algorithm~2, zero suspicious contiguous zero-runs) and a dynamic
check pushing a generic basis vector through one real full round at both
BN254 and KoalaBear scale. We additionally confirmed this is not merely a
paper construction: reading Plonky3's production
\texttt{mds\_light\_permutation} code directly shows it uses the
\emph{unswapped}, vulnerable order ($P_{t/4}\otimes M_4$) in a real,
actively used STARK proving stack, while \texttt{circomlib}'s
non-tensor-structured reference matrix is clean under the same scan.

% =========================================================================
\section{Replicating and Defeating Jo's Full-Round Collision}
\label{sec:midpoint-reset}
% =========================================================================

Jo's Discussion section proposes, without building or testing it, exactly
the binding countermeasure of \S\ref{sec:binding}: \emph{``derive the
complete parameter set non-adaptively\dots by binding the round constants
to a canonical encoding of the matrix.''} We built and validated this
claim.

\subsection{The attack mechanism}

Track two executions' midpoint $m_r = (x^+_r+x^-_r)/2$ and half-difference
$d_r=(x^+_r-x^-_r)/2$. One prescribed matrix image per two-round block,

\[
  M u_r = -c^{(r+1)},
\]

cancels the midpoint against the next round constant; since the S-box is
odd ($S(-z)=-S(z)$ for odd $\alpha$), the following S-box's output
midpoint is exactly zero and stays zero for the rest of the permutation
(``midpoint reset''). Two further reusable images, $Me=v$ and $Mw=e$ on a
coordinate $e$ passive in every partial round, keep the half-difference
confined to the line $\Fp\cdot e$ through every full \emph{and} partial
round. Collecting these prescribed images as columns of $U$ (sources,
computable from the public constants alone --- not from $M$) and $V$
(targets), $M = VU^{-1}$ is the unique matrix realizing the whole
trajectory. A correctly closed scalar trajectory (final difference equal
to the negative of the initial difference) yields a genuine feed-forward
collision in an $M$ that is MDS and passes every subspace-trail check in
the literature.

\subsection{Replication}

We implemented the construction generically (any $t$, field, round
constants), with one deliberate scope reduction relative to Jo's own
$(t,\alpha,R_F,R_P)=(16,3,8,20)$ KoalaBear instance: $R_F=4$ (one
full-round block per side of the partial-round segment, instead of two).
This collapses Jo's harder two-block cubic entry/exit solve to a
closed-form single step: choosing the entry block's free scalar $s_0=0$
gives $b_0 = \sqrt[3]{a}$ directly (well-defined since $\gcd(3,p-1)=1$,
the same condition $\alpha=3$ already requires), and the block's output
lands at exactly $a$ by construction. The one remaining free scalar (the
closing amplitude $a$) is found by brute-force search over a small toy
prime field, in the same spirit as the project's existing policy of
replicating a mechanism at reduced scale rather than its literal headline
instance.

\paragraph{Result.} Against unprotected (unbound) constants, at
$t=6,\alpha=3,R_F=4,R_P=4$ over a toy field $p=10037$, a closing amplitude
$a^{*}=177$ yields a genuine, independently-verified feed-forward
collision:
\[
  H(x^+) = H(x^-) = (0,\,3518,\,0,\,0,\,0,\,0), \qquad x^+ \neq x^-,
\]
confirmed by executing the actual permutation forward with the solved $M$
(not merely the scalar bookkeeping). A real implementation bug was caught
in the process: an initial version computed $M = U^{-1}V$ instead of
$M=VU^{-1}$ (matrix multiplication is non-commutative), silently solving
a different linear system; caught by a direct $M\,U \stackrel{?}{=} V$
sanity check, now retained in the code path.

\subsection{Does binding close the gap?}

We generated a matrix $M$ the honest way (NUMS,
\S\ref{sec:lightweight}'s $\MI$ family), derived constants
$c = \mathrm{Hash}(\mathtt{seed}\,\|\,\mathrm{params}\,\|\,M\,\|\,M)$ via
the binding of \S\ref{sec:binding} --- constants fixed \emph{after and
from} $M$, the reverse of the order the attack needs --- and then scanned
\emph{every} nonzero $a$ in the toy field for whether this already-fixed
$M$ happens to satisfy the attack's required image constraints
$M\,U(c,a) = V(c,a)$ for any of them.

\begin{table}[h!]
\centering
\caption{Binding-immunity test: consistent closing amplitudes found out
of $p-1$ candidates scanned, across independent trials.}
\label{tab:immunity}
\begin{tabular}{@{}lllcc@{}}
\toprule
Trial set & Field $p$ & $\MI$ family & Candidates scanned & Consistent $a$ found \\
\midrule
5 independent seeds & 10007 & \texttt{poseidon2\_sparse} & 10006 each & 0 \\
Cross-family & 10007 & \texttt{low\_weight\_circulant} & 10006 & 0 \\
Cross-family & 10007 & \texttt{irreducible\_filtered} & 10006 & 0 \\
Cross-field & 10067 & \texttt{low\_weight\_circulant} & 10066 & 0 \\
Cross-field & 10067 & \texttt{irreducible\_filtered} & 10066 & 0 \\
\bottomrule
\end{tabular}
\end{table}

Zero consistent amplitudes were found in \emph{every} trial (nine
independent trials total, Table~\ref{tab:immunity}). We state the result
at the confidence it supports: this is not a formal theorem that binding
defeats every instance of this attack class, but direct empirical evidence
that honestly NUMS-generated, bound $(M,c)$ pairs do not happen to satisfy
the attack's highly constrained algebraic relation --- consistent with the
heuristic that an $M$ uncorrelated with the attack's image constraints by
construction satisfies them only with probability on the order of the
field size, and with Jo's own qualitative argument that binding removes
the ``choose $M$ after $c$'' degree of freedom the construction depends
on. A formal fixed-point/preimage-hardness argument, and replication at
Jo's literal $R_F=8$/KoalaBear scale (which needs a general cubic-root
solver rather than our closed-form shortcut), are left as future work.

% =========================================================================
\section{Lightweight Instantiation}
\label{sec:lightweight}
% =========================================================================

\subsection{Candidate generation and cost model}

Candidate $\MI$ matrices are generated via three NUMS families:
\texttt{poseidon2\_sparse} (near-identity off-diagonal, hash-derived
diagonal --- the real Poseidon2 internal-matrix shape), a low-weight
circulant, and a dense matrix rejection-sampled for GRS21 irreducibility.
We separate two cost measures that the literature sometimes conflates:
\[
  \mathrm{cost}_{\mathrm{R1CS}}(t,R_F,R_P,\alpha) = R_F\cdot t \cdot
  c_\alpha + R_P \cdot c_\alpha,
  \qquad
  \mathrm{mul}(M) = |\{(i,j): M_{ij}\notin\{0,1\}\}|,
\]
where $c_\alpha$ is the number of multiplication constraints for
$x^{\alpha}$ via square-and-multiply ($c_5=3$, $c_3=2$). In a standard
R1CS encoding, multiplying a wire by a public constant is a free linear
combination; $\mathrm{cost}_{\mathrm{R1CS}}$ is S-box-dominated and
independent of the matrix, while $\mathrm{mul}(M)$ is the real cost
differentiator for native/Plonkish backends.

\begin{table}[h!]
\centering
\caption{Candidate comparison at $t=3,\alpha=5,R_F=8,R_P=57$ (BN254).}
\label{tab:candidates}
\begin{tabular}{@{}lccc@{}}
\toprule
Family & R1CS non-linear & R1CS linear & $\mathrm{mul}(\MI)$ (per round) \\
\midrule
\texttt{poseidon2\_sparse} & 243 & 276 & 195 \\
\texttt{low\_weight\_circulant} & 243 & 276 & 390 \\
\texttt{irreducible\_filtered} & 243 & 276 & 585 \\
\bottomrule
\end{tabular}
\end{table}

As Table~\ref{tab:candidates} shows, R1CS constraint count is
\textbf{identical} across all three candidates --- confirmed by an
independent \texttt{circom} compile, not merely the formula --- because a
circuit row is allocated per matrix \emph{row}, not per nonzero entry.
Matrix sparsity is a real win for native/Plonkish multiplicative
complexity and has \textbf{zero} measurable effect on R1CS/Groth16 cost.
\texttt{poseidon2\_sparse} is the clear winner on the one metric that
differs.

\subsection{Cross-validation against the production reference}

Compiling our NUMS-generated circuit and \texttt{iden3/circomlib}'s
production \texttt{Poseidon(2)} template at identical parameters
($t=3,R_F=8,R_P=57,\alpha=5$) gives \textbf{byte-identical} non-linear
constraint counts (243 for both) and near-identical linear counts (276 vs.
274, the 2-constraint difference being cosmetic hash-wrapper exposure).
A full Groth16 pipeline (Powers-of-Tau, setup, prove, verify) produced
valid, independently-verified proofs for both circuits with
indistinguishable wall-clock timing at this circuit size.

\subsection{The field-size lever}

Generalizing the field (BN254 $\to$ KoalaBear, $p=2^{31}-2^{24}+1$) is a
real, verified lightweight lever, with one correction found in the
process. A first attempt reused $R_F=8$ from the public $t=16$ KoalaBear
instance without re-deriving it for $t=3$; GKR+21's own statistical
security bound (Eq.~2, which scales with a $(t+1)$ factor) requires
$R_F=12$ at $t=3$, not $8$. With the corrected round count,
$t=3,\alpha=3,R_F=12,R_P=20$ over KoalaBear gives $112$ non-linear $+ 152$
linear $= 264$ total constraints, versus $512$--$519$ at the BN254
baseline --- a real $\approx 1.9\times$ reduction, confirmed by a
\texttt{circom --prime goldilocks} compile. This is a deployment decision
(it requires the surrounding system to not be anchored to BN254/Groth16),
not an algorithmic improvement, and is clearly weaker than the naive
$2.4\times$ an uncorrected round count would have claimed.

\subsection{Tensor-structuring $\Mfull$ too}

Every candidate above optimizes only $\MI$; $\Mfull$ was always a dense
NUMS matrix. Applying the safe-order tensor structure
(\S\ref{sec:tensor}) to $\Mfull$ as well, at $t=16$ (Vitto's own instance
width), gives identical R1CS cost (confirmed by compile: $296$ non-linear
$+ 596$ linear $= 892$, byte-identical to the all-dense control) but a
substantial real multiplicative-complexity win: a naive entry count gives
$184$ nonzero-non-one entries, but the real fast block algorithm (apply
the dense $4\times4$ block $M_4$ to each of $t/4$ blocks, then a
pure-addition outer mix) costs only $4\times 10=40$ --- a $4.6\times$
difference the naive count misses entirely. Recomputing the full
permutation's multiplicative complexity with the honest fast-algorithm
number gives $8\times 40 + 20\times 16 = 640$ versus $8\times 256 +
20\times 16 = 2368$ for the all-dense design: a \textbf{73\% reduction},
not the 24.3\% the naive entry-count comparison first suggested.

% =========================================================================
\section{Exhaustive Validation of the Final Candidate}
\label{sec:battery}
% =========================================================================

The final candidate ($t=16,\alpha=3,R_F=8,R_P=20$, KoalaBear,
$\MI=\texttt{poseidon2\_sparse}$, $\Mfull = M_4\otimes P_4$ in the safe
order, constants bound to both matrices) was run through every available
check (Table~\ref{tab:battery}).

\begin{table}[h!]
\centering
\caption{Exhaustive battery against the final candidate.}
\label{tab:battery}
\small
\begin{tabular}{@{}p{0.05\linewidth}p{0.42\linewidth}p{0.45\linewidth}@{}}
\toprule
\# & Check & Result \\
\midrule
1 & Statistical round sufficiency (GKR+21 Eq.~2) & $\min R_F = 8$ (bound 6, margin 2); candidate $R_F=8$ --- sufficient, exactly at the boundary \\
2 & MDS / branch number of $\MI$ & $B(\MI)=17=t+1$ --- MDS, no slack \\
3 & Algorithm 1 --- infinite invariant subspace in $\MI$ & \texttt{SECURE} \\
4 & Algorithm 2 --- static structural scan of $\Mfull$ & flags one subspace; investigated and confirmed a \textbf{false positive} (resolved by item 7) \\
5 & Finite-trail dimension profile of $\MI$ & $\dim K_r = t-r$ for every $r=0,\dots,15$ --- maximal rank growth \\
6 & Symbolic freeze, adversarial $d\in K_4(\MI)$ & entering degrees $[0,0,0,0,1,3]$ --- frozen exactly 4 rounds, jumps at round 4, matching $\dim K_4=12$ exactly \\
7 & No-accidental-landing, generic direction through $\Mfull$ & $[1,3,9]$, matching naive $\alpha^{\text{round}}$ growth exactly --- no accidental trail \\
8 & Real multiplicative cost & $640$ (cross-confirmed against this exact candidate) \\
\bottomrule
\end{tabular}
\end{table}

\paragraph{The item-4 false positive, resolved.} The static scan flags
the \emph{safe}-order $\bar M_\varepsilon = M_4\otimes P_4$ at KoalaBear
scale, where it was clean at BN254 --- not a contradiction, since which
irreducible factors $\chi(\bar M_\varepsilon)$ splits into over $\Fp$ is
field-dependent. Pushing a \emph{generic} (non-$\pm1$) basis vector of the
flagged subspace through one real full round shows it reaches full weight
(no negation-compatible structure survives), confirming the flag is a
genuine-but-inert invariant subspace, not an exploitable one. We state
the resulting methodological lesson explicitly: the static scan is
necessary-looking but not sufficient, and must always be paired with a
one-round dynamic survival check before a flag is reported as a finding.

\subsection{Comparison against the two mechanism papers}

\begin{table}[h!]
\centering
\caption{Precondition each attack needs, and why this candidate's design removes it.}
\label{tab:comparison}
\small
\begin{tabular}{@{}p{0.28\linewidth}p{0.32\linewidth}p{0.32\linewidth}@{}}
\toprule
Attack & Precondition needed & This candidate's defense \\
\midrule
Full-round absorption (Vitto) & Matrix chosen knowing the constants (or vice versa) & Both matrices bound into one NUMS transcript before either is derived (\S\ref{sec:binding}) \\
Round-skipping (Merz \& Garc\'ia) & Unswapped Kronecker order & Swapped order from the start (\S\ref{sec:tensor}), verified at both BN254 and KoalaBear \\
\bottomrule
\end{tabular}
\end{table}

Neither paper's literal headline numbers (Vitto's $3^{10}$ vs.\ $3^{28}$
collapse; Merz and Rodr\'iguez Garc\'ia's concrete query complexity) are
reproduced bit-for-bit here --- that needs their full adversarial
multi-group construction, a separate undertaking from checking whether
this candidate resists the underlying mechanism, which it does.

% =========================================================================
\section{Real-World Audit}
% =========================================================================

We ran Algorithms~1 and~2 against three production parameter sets not
generated by this project:

\begin{itemize}[leftmargin=1.4em]
  \item \textbf{\texttt{circomlib}}: the real $t=3$ matrix from the
  widely-deployed reference circuit is clean (MDS, Algorithm~1
  \texttt{SECURE}, Algorithm~2 reports no suspicious subspace).
  \item \textbf{Plonky3}: round numbers match our own
  statistically-derived formula and the current SoK's recomputed table
  exactly; the external matrix uses the \emph{vulnerable} (unswapped)
  Kronecker order, confirmed three independent ways (direct code reading,
  reproducing the matrix and exhibiting the attack shape, and our
  Kronecker-agnostic Algorithm~2 flagging it). Practical exploitability
  depends on the sponge rate/capacity split at each call site, which we
  did not chase down.
  \item \textbf{\texttt{arkworks-rs}}: ships no fixed matrix at all ---
  every consumer calls a Cauchy-matrix generator whose own source comments
  claim three well-formedness checks that are \textbf{asserted in a
  comment with zero corresponding code}. This is a documentation/
  implementation gap, not a demonstrated break of any specific deployed
  instance, but it confirms that a widely-depended-on parameter generator
  provides no automated protection against generating a GRS21- or
  Vitto-vulnerable matrix for an arbitrary $(t,\text{field})$ choice.
\end{itemize}

% =========================================================================
\section{Positioning Against the Efficiency State of the Art}
% =========================================================================

A full read of the most comprehensive available survey, ``SoK:
Zero-Knowledge Friendly Hash Functions over Prime
Fields''~\cite{sok2026} (134 pages, read in full for this paper
after earlier attempts were blocked by access restrictions), gives real,
already-computed R1CS/Groth16 constraint counts we cite directly rather
than re-deriving (Table~\ref{tab:sok}).
\label{sec:sok}

\begin{table}[h!]
\centering
\caption{R1CS constraint counts, single permutation, Gnark/Groth16,
BLS12-381, $t=4$ (SoK~\cite{sok2026}, Table~23).}
\label{tab:sok}
\begin{tabular}{@{}lc@{}}
\toprule
Construction & R1CS constraints \\
\midrule
Griffin & 224 \\
Anemoi~\cite{anemoi2022} & 244 \\
Arion & 254 \\
Poseidon & 268 \\
Poseidon2 & 268 \\
Rescue-Prime & 368 \\
Skyscraper & 1425 \\
Polocolo~\cite{polocolo2025} & 1553 \\
\bottomrule
\end{tabular}
\end{table}

Two findings from this read directly inform this paper's scope. First,
the SoK's attack-cryptanalysis section (\S4, organized as a cross-cutting
technique toolbox rather than one checkable criterion) mentions both
Vitto's and Jo's attacks in a single combined sentence, with no mechanism
detail, no countermeasure, and no round-number consequence drawn ---
confirming that the unified criterion and the binding validation of
\S\ref{sec:criterion}--\S\ref{sec:midpoint-reset} are not duplicated
anywhere in the existing literature. Second, the SoK confirms
\emph{``no instances require updating under the latest attacks''} for
Poseidon/Poseidon2's own round-number formula, an independent validation
of the statistical bound we use in \S\ref{sec:battery}, item~1.

% =========================================================================
\section{Discussion: What Did Not Work, and What Remains Open}
\label{sec:discussion}
\label{sec:levers}
% =========================================================================

\subsection{Negative results, reported honestly}

\begin{itemize}[leftmargin=1.4em]
  \item \textbf{Lowering $\alpha$.} At $t=3$, reducing $\alpha$ from $5$
  to $3$ requires $R_P$ to grow from $56$ to $83$; the extra rounds'
  linear-constraint overhead more than cancels the per-S-box saving
  ($594$ vs.\ $512$ total constraints, confirmed by compile). This lever
  makes the design \emph{heavier}, not lighter, in R1CS.
  \item \textbf{Rotating the active S-box lane.} Generalizing which
  coordinate carries the partial-round S-box offers no measured benefit
  for matrices already at the finite-trail dimension-decay ceiling (our
  actual candidates), and can \emph{introduce} a new infinite trail in a
  matrix that was previously secure under the fixed-position design.
  \item \textbf{Two tricks imported from GMiMCHash2.} The $\alpha=2^k$
  S-box does not port: GMiMC is a Feistel construction where the round
  function need not be a bijection, while Poseidon's SPN structure
  requires $\gcd(\alpha,p-1)=1$ for invertibility, which no power of two
  satisfies over an odd-characteristic field. The $M_{IO}$ input/output
  matrix targets analytical tractability of a designer's own security
  proof, not a round-count or constraint-count reduction, and moves none
  of our cost numbers.
  \item \textbf{A naive lookup-hybrid model.} Modeling a ``lookup round''
  as contributing near-maximal algebraic degree in a multiplicative
  degree-bar tracker produces a vacuous result: it reports security for
  replacing $56$ of $57$ partial rounds with a single lookup round,
  because the model has no way to represent an attacker simply ignoring
  the lookup round and attacking the short algebraic remainder directly.
  A real answer needs either a close reading of how lookup-based designs
  formally argue hybrid security, or a lookup-capable backend (Halo2,
  Plonky2), neither pursued here.
\end{itemize}

\subsection{Open problems}

Threats~4 and~5 of Table~\ref{tab:threats} --- S-box skipping (Bhati,
Tariq, and Ashur, eprint 2026/1692) and bivariate resultant attacks
against Poseidon's own linear layer (eprint 2026/1905) --- remain open.
Both are cited only in passing, without mechanism analysis, in the most
comprehensive survey available~\cite{sok2026}; closing either with the
same Phase-0 replication discipline used for
\S\ref{sec:threats}--\S\ref{sec:midpoint-reset} is the natural next step
for this line of work.

% =========================================================================
\section{Conclusion}
% =========================================================================

MDS is necessary for Poseidon's partial-round linear layer, but three
independent 2026 results show it is not sufficient: properly-MDS matrices
that pass every published subspace-trail check can still be broken via
finite-trail absorption, tensor-structured round skipping, or an
adaptively chosen matrix that programs an exact collision trajectory. We
gave a unified structural criterion covering the first two mechanisms, a
single generation-order countermeasure (binding round constants to both
linear layers before either is used) that structurally removes the
precondition all three mechanisms need, and validated that countermeasure
empirically against all three attacks --- including, to our knowledge,
the first replication and defense validation of Jo's full-round collision
reported anywhere. We combined this with a lightweight KoalaBear
instantiation, cross-validated against a production reference
implementation and an exhaustive eight-point falsification battery, and
positioned it against the current efficiency state of the art. S-box
skipping and bivariate resultant attacks against Poseidon's own linear
layer remain open problems for future work.

% =========================================================================
\begin{thebibliography}{99}

\bibitem{gkr21}
L.~Grassi, D.~Khovratovich, C.~Rechberger, A.~Roy, M.~Schofnegger,
\emph{Poseidon: A New Hash Function for Zero-Knowledge Proof Systems},
USENIX Security 2021.

\bibitem{gks23}
L.~Grassi, D.~Khovratovich, M.~Schofnegger,
\emph{Poseidon2: A Faster Version of the Poseidon Hash Function},
AFRICACRYPT 2023.

\bibitem{hades20}
L.~Grassi, R.~L\"uftenegger, C.~Rechberger, D.~Rotaru, M.~Schofnegger,
\emph{On a Generalization of Substitution-Permutation Networks: The HADES
Design Strategy}, EUROCRYPT 2020.

\bibitem{grs21}
L.~Grassi, C.~Rechberger, M.~Schofnegger,
\emph{Proving Resistance Against Infinitely Long Subspace Trails: How to
Choose the Linear Layer}, IACR ToSC 2021(2):314--352.

\bibitem{vitto2026}
G.~Vitto, \emph{Slipway: Accessing Finite Subspace Trails in Poseidon},
IACR ePrint 2026/1579.

\bibitem{merzgarcia2026}
S.-P.~Merz, \`A.~Rodr\'iguez Garc\'ia,
\emph{Skipping Class: Algebraic Attacks Exploiting Weak Matrices and
Operation Modes of Poseidon2(b)}, IACR ePrint 2026/306.

\bibitem{jo2026}
S.~Jo, \emph{Midpoint Reset: A Full-Round Poseidon Collision from an
Adaptively Chosen MDS Matrix}, IACR ePrint 2026/1760.

\bibitem{backdoor2026}
\emph{Nothing Up My Matrix: Kleptographic Backdoors in ZK-friendly Hash
Functions}, IACR ePrint 2026/1901.

\bibitem{sboxskip2026}
A.~S.~Bhati, S.~Tariq, T.~Ashur,
\emph{From Round Skipping to S-Box Skipping: Attacking Poseidon's Partial
Layer via Subspace Restriction}, IACR ePrint 2026/1692.

\bibitem{resultant2026}
A.~Bak, M.~Hostettler,
\emph{A Better Bivariate Resultant Attack on Round-Reduced Poseidon},
IACR ePrint 2026/1905.

\bibitem{sok2026}
C.~Bouvier, L.~Grassi, K.~Koschatko, C.~Rechberger, F.~Schmid,
M.~J.~Steiner, Z.~Wu, H.~Yan,
\emph{SoK: Zero-Knowledge Friendly Hash Functions over Prime Fields},
IACR ePrint 2026/1931.

\bibitem{anemoi2022}
C.~Bouvier, P.~Briaud, P.~Chaidos, L.~Perrin, et al.,
\emph{New Design Techniques for Efficient Arithmetization-Oriented Hash
Functions: Anemoi Permutations and Jive Compression Mode},
IACR ePrint 2022/840.

\bibitem{polocolo2025}
J.~Ha, S.~Hwang, J.~Lee, S.~Park, M.~Son,
\emph{Polocolo: A ZK-Friendly Hash Function Based on S-boxes Using Power
Residues}, EUROCRYPT 2025.

\end{thebibliography}

\end{document}
